% Work-MAT | 2. rocnik | pripravovana hodina 52
% Návaznost: uživatelské připravy.zip (hodiny 01–20), autorská učebnice na GitHubu.
% Realisticky.cz: didaktická inspirace v posloupnosti a typech úloh; zadání jsou vlastní.
\documentclass[11pt,a4paper]{article}
\input{../../sablony/pripravy/priprava.sty}
\usepackage{tikz}
\nastavPripravu{2. ročník}{Ověřování znalostí}{52}{Pololetní písemná práce P2}
\begin{document}
\blok{Motivace a návaznost}
Pololetní práce má vycházet z aktuálního rozsahu: exponenciální a logaritmické rovnice, úvod do goniometrie. Zde je metodické rozvržení a analogický rozbor, nikoli testové varianty.
\poznamka{Orientační nabídka pro 45--55 minut: krátké zopakování a motivace, nový výklad, tři společné vzorové úlohy, samostatná práce. Vyber podle tempa třídy; zbývající zadání slouží jako rezerva.}
\blok{Výklad}
Ve chvíli rozboru odděl chybu ve vzorci od chyby v podmínce či znaménku hodnoty funkce.
\blok{Společné příklady}
\begin{prikladReseniBox}[title={Příklad 1}]
Vyřeš, resp. urči: $3^{x+1}+3^x=108$.
\tcblower
\[\begin{aligned}4\cdot3^x=108 \\ x=3\end{aligned}\]
\end{prikladReseniBox}
\begin{prikladReseniBox}[title={Příklad 2}]
Vyřeš, resp. urči: $\sin150^\circ,\quad\cos150^\circ$.
\tcblower
\[\begin{aligned}\sin150^\circ=\frac{1}{2} \\ \cos150^\circ=-\frac{\sqrt3}{2}\end{aligned}\]
\end{prikladReseniBox}
\poznamka{Požaduj schválené testové varianty; hodnoticí kritéria a konkrétní podmínky vždy ověř podle školy.}
\blok{Samostatné procvičování}
\begin{zadaniReseniBox}[title={Samostatná práce 1}]
\begin{enumerate}[label=\alph*)]
\item $\operatorname{tg}45^\circ$.
\item $\sin0^\circ$.
\item $\operatorname{cotg}135^\circ$.
\item $\sin60^\circ$.
\item $\sin30^\circ$.
\item $\cos60^\circ$.
\end{enumerate}
\tcblower
\begin{enumerate}[label=\alph*)]
\item $1$.
\item $0$.
\item $-1$.
\item $\frac{\sqrt3}{2}$.
\item $\frac{1}{2}$.
\item $\frac{1}{2}$.
\end{enumerate}
\end{zadaniReseniBox}
\blok{Dobrovolné rozšíření}
\begin{zadaniReseniBox}[title={Rozšíření 1}]
\begin{enumerate}[label=\alph*)]
\item Navrhni obdobnou úlohu s jiným úhlem a zkontroluj výsledek graficky..
\end{enumerate}
\tcblower
\begin{enumerate}[label=\alph*)]
\item $Srovnej s jednotkovou kružnicí nebo grafem.$.
\end{enumerate}
\end{zadaniReseniBox}
\blok{Úkol na doma}
Každou hodnotu zdůvodni náčrtkem jednotkové kružnice nebo orientovaným úhlem; u přesných hodnot neuváděj jen desetinné aproximace.
\begin{zadaniReseniBox}[title={Domácí úkol}]
\begin{enumerate}[label=\alph*)]
\item $\sin45^\circ$.
\item $\cos45^\circ$.
\item $\sin150^\circ$.
\item $\cos150^\circ$.
\item $\operatorname{tg}135^\circ$.
\item $\operatorname{cotg}45^\circ$.
\item $\sin210^\circ$.
\item $\cos240^\circ$.
\end{enumerate}
\tcblower
\begin{enumerate}[label=\alph*)]
\item $\frac{\sqrt2}{2}$.
\item $\frac{\sqrt2}{2}$.
\item $\frac{1}{2}$.
\item $-\frac{\sqrt3}{2}$.
\item $-1$.
\item $1$.
\item $-\frac{1}{2}$.
\item $-\frac{1}{2}$.
\end{enumerate}
\end{zadaniReseniBox}
\poznamka{Domácí úkol obsahuje 8 položek; předpokládaná domácí práce přibližně 15--30 minut. Po větších blocích jsou vědomě promíchané dříve probrané typy.}
\end{document}
